% !TEX TS-program = pdflatex
% !TeX program = pdflatex
% !TEX encoding = UTF-8
% !TEX spellcheck = en_US

\documentclass{beamer}

\usepackage{../extra/beamer/karimnlp}


%\usepackage{fullpage}
%\usepackage[left=2.8cm,right=2.2cm,top=2 cm,bottom=2 cm]{geometry}
\setbeamersize{text margin left=10pt,text margin right=10pt}



%\usepackage[T1]{fontenc}
\usepackage{textcomp}
\usepackage[utf8]{inputenc}
%\usepackage[french,english]{babel}
\usepackage{arabtex}
\usepackage{txfonts}
\usepackage{tipa}
\usepackage[]{graphicx}
\usepackage{multirow}
\usepackage{hyperref}
\usepackage{colortbl}
\usepackage{tabularray}
\usepackage{listingsutf8}
\usepackage{wrapfig}
\usepackage{multicol}
\usepackage[export]{adjustbox} %for images in table, also for frame
\usepackage[many]{tcolorbox}
\usepackage{wasysym}
\usepackage[lined]{algorithm2e}
\usepackage{alltt} %verbatim with commands
\usepackage{longtable}
\usepackage{tabu}

\usepackage{CJKutf8}

\usepackage{amsmath,amssymb, amsfonts} 

%\usepackage{newtxtext,newtxmath}


%\usepackage{acolor}


\definecolor{lightblue}{HTML}{D0D2FF}
\definecolor{lightyellow}{HTML}{FFFFAA}
\definecolor{darkblue}{HTML}{0000BB}
\definecolor{olivegreen}{HTML}{006600}
\definecolor{darkgreen}{HTML}{008B45} %009B55
\definecolor{violet}{HTML}{6600CC}
\definecolor{deeppink}{HTML}{FF1493}
\definecolor{orangey}{HTML}{FFBB00}


\makeatletter
\def\beamer@calltheme#1#2#3{%
	\def\beamer@themelist{#2}
	\@for\beamer@themename:=\beamer@themelist\do
	{\IfFileExists{\beamer@themelocation/#3\beamer@themename.sty}{%
	\usepackage[{#1}]{\beamer@themelocation/#3\beamer@themename}}{\usepackage[{#1}]{#3\beamer@themename}}
		
	}%
}

\def\usefolder#1{
	\def\beamer@themelocation{#1}
}
\def\beamer@themelocation{}
\makeatother

\usefolder{../../extra/beamer-themes}

\usetheme{Karim} % Antibes Boadilla Warsaw Karim

\beamertemplatenavigationsymbolsempty

\let\oldcite\cite
\renewcommand{\cite}[1]{{\bfseries\color{blue}\oldcite{#1}}}

\tcbuselibrary{listings}


%\renewcommand{\baselinestretch}{1.5}

\def\supit#1{\raisebox{0.8ex}{\small\it #1}\hspace{0.05em}}

\AtBeginSection{%
	\begin{frame}
		\sectionpage
	\end{frame}
}

\newcommand{\rottext}[2]{%
	\rotatebox{90}{%
	\begin{minipage}{#1}%
		\raggedleft#2%
	\end{minipage}%
	}%
}


\institute{ %
Laboratoire de la Communication dans les Systèmes Informatiques (LCSI)
\par
École  nationale Supérieure d'Informatique (ESI), Algiers, Algeria
}
\author[Abdelkrime Aries (2023/2024)] %
{Abdelkrime Aries}
%\titlegraphic{\includegraphics[height=1cm]{../img/esi-logo.png}%\hspace*{4.75cm}~


\date{Academic year: 2023/2024} %\today

\titlegraphic{%
	\includegraphics[height=1cm]{../../img/misc/esi-logo.png}%
	\hspace{2cm}%
	\includegraphics[height=1cm]{../../img/misc/lcsi-logo.png}%
	\hspace{2cm}%
	\includegraphics[height=1.5cm]{../../img/misc/esi.nlp-logo.png}%
}

%\setbeamertemplate{headline}{}

\newcommand{\kurl}[1]{{\scriptsize\bfseries\color{orangey}\url{#1}}}

\newcommand{\keyword}[1]{\textcolor{red}{\bfseries\itshape #1}}
\newcommand{\expword}[1]{\textcolor{olivegreen}{#1}}
\newcommand{\optword}[1]{\textcolor{violet}{\bfseries #1}}

\makeatletter
\newcommand\mysphere{%
	\parbox[t]{10pt}{\raisebox{0.2pt}{\beamer@usesphere{item projected}{bigsphere}}}}
\makeatother

%\let\oldtabular\tabular
%\let\endoldtabular\endtabular
%\renewenvironment{tabular}{\rowcolors{2}{white}{lightblue}\oldtabular\rowcolor{blue}}{\endoldtabular}


%\NoAutoSpacing %french autospacing after ":"

\def\graphpath{}

\newcommand{\changegraphpath}[1]{\def\graphpath{#1}}


\newcommand{\vgraphpage}[2][.84\textheight]{%
%	\begin{center}%
		\includegraphics[height=#1]{\graphpath #2}%
%	\end{center}%
}

\newcommand{\hgraphpage}[2][\textwidth]{%
%	\begin{center}%
		\includegraphics[width=#1]{\graphpath #2}%
%	\end{center}%
}

\newcommand{\graphpage}[2][]{%
	\includegraphics[#1]{\graphpath #2}%
}

\bibliographystyle{apalike}

\newcommand{\insertbibliography}[2]{
	\appendix
	\section*{Bibliography}
	\nocite{#2}
%	\makeatletter % to change template
%	\setbeamertemplate{headline}[default] % not mandatory, but I though it was better to set it blank
%	\def\beamer@entrycode{\vspace*{-\headheight}} % here is the part we are interested in :)
%	\makeatother
	\begin{multicols*}{2}[\frametitle{\insertsection} \usebeamertemplate*{frametitle}]%\usebeamertemplate*{frametitle}\frametitle{Références}
		\tiny
		\bibliography{#1}
	\end{multicols*}
}

\definecolor{my-grey}{RGB}{233, 233, 233}

\newcommand{\insertlicence}{
	\begin{frame}[plain]
	\frametitle{License: CC-BY 4.0}

	\begin{tcolorbox}[colback=cyan,
		colframe=cyan,  
		arc=0pt,outer arc=0pt,
		valign=top, 
		halign=center,
		width=\textwidth]
		
		\includegraphics[width=.5cm]{../../img/licence/cc_icon_white_x2.png}
		\includegraphics[width=.5cm]{../../img/licence/attribution_icon_white_x2.png}
		
		\color{white}
		\bfseries Attribution 4.0 International (CC BY 4.0) \\
		\tiny \url{https://creativecommons.org/licenses/by/4.0/deed.en}
		
	\end{tcolorbox}\vspace{-.5cm}
	\begin{tcolorbox}[colback=my-grey,
		colframe=my-grey,  
		center, arc=0pt,outer arc=0pt,
		valign=top, 
		halign=left,
		width=\textwidth]
		
		\tiny
		
		\begin{center}
			\bfseries\Large
			You are free to:
		\end{center}
		
		\begin{minipage}{0.83\textwidth}
			\begin{itemize}
				\item[] \textbf{Share} — copy and redistribute the material in any medium or format
				\item[] \textbf{Adapt} — remix, transform, and build upon the material
				for any purpose, even commercially.
			\end{itemize}
		
			\vspace{12pt}
		
			\textit{The licensor cannot revoke these freedoms as long as you follow the license terms.}
		\end{minipage}
		\begin{minipage}{0.15\textwidth}
			\includegraphics[width=\textwidth]{../../img/licence/FreeCulturalWorks_seal_x2.jpg}
		\end{minipage}
	
		
		\begin{center}
			\bfseries\Large
			Under the following terms:
		\end{center}
		
		\begin{itemize}
			\item[] \textbf{Attribution} — You must give appropriate credit, provide a link to the license, and indicate if changes were made. You may do so in any reasonable manner, but not in any way that suggests the licensor endorses you or your use.
			\item[] \textbf{No additional restrictions} — You may not apply legal terms or technological measures that legally restrict others from doing anything the license permits.
		\end{itemize}
		
	\end{tcolorbox}
	
%	\begin{center}
%		\bfseries Attribution 4.0 International (CC BY 4.0)
%		\url{https://creativecommons.org/licenses/by/4.0/deed.fr}
%	\end{center}

%	\tiny
%
%	Vous êtes autorisé à : 
%	\begin{itemize}
%		\item \textbf{Partager} — copier, distribuer et communiquer le matériel par tous moyens et sous tous formats
%		\item \textbf{Adapter} — remixer, transformer et créer à partir du matériel
%	\end{itemize}
%	
%	Selon les conditions suivantes : 
%	\begin{itemize}
%		\item \textbf{Attribution} — Vous devez créditer l'Œuvre, intégrer un lien vers la licence et indiquer si des modifications ont été effectuées à l'Oeuvre. Vous devez indiquer ces informations par tous les moyens raisonnables, sans toutefois suggérer que l'Offrant vous soutient ou soutient la façon dont vous avez utilisé son Oeuvre.
%		\item \textbf{Pas d'Utilisation Commerciale} — Vous n'êtes pas autorisé à faire un usage commercial de cette Oeuvre, tout ou partie du matériel la composant. 
%		\item \textbf{Pas de restrictions complémentaires} — Vous n'êtes pas autorisé à appliquer des conditions légales ou des mesures techniques qui restreindraient légalement autrui à utiliser l'Oeuvre dans les conditions décrites par la licence.
%	\end{itemize}

	\end{frame}
}

\settowidth{\leftmargini}{\usebeamertemplate{itemize item}}
\addtolength{\leftmargini}{\labelsep}

\AtBeginDocument{
	\newcolumntype{L}[2]{>{\vbox to #2\bgroup\vfill\flushleft}p{#1}<{\egroup}} 
	
	\begin{frame}[plain]
		\maketitle
	\end{frame}

	\insertlicence
}


% needs etoolbox; to break links after -
\appto\UrlBreaks{\do\-}


\makeatletter
\newcommand{\xRightarrow}[2][]{\ext@arrow 0359\Rightarrowfill@{#1}{#2}}
\makeatother


\usefonttheme{structurebold}
%\usefonttheme{professionalfonts}


% figure caption
\setlength\abovecaptionskip{2pt}

% equations
\AtBeginDocument{
\abovedisplayskip=0pt
\abovedisplayshortskip=0pt
\belowdisplayskip=0pt
\belowdisplayshortskip=0pt
}


%\usepackage{stix}

\subtitle[Presentation]{Course ``Curse" presentation} 

\begin{document}
 

\begin{frame}
\frametitle{Prerequisites}

\begin{itemize}
	\item Basic knowledge of languages (vocabulary, grammar, etc.);
	\item Data structures (1CP-ALSDD);
	\item Probability and statistics \\(2CP-PRST1, 2CP-PRST2);
	\item Algebra and calculus \\(2CP-ANAL3, 2CP-ALG3, 2CP-ANAL4);
	\item Language theory (1CS-THP);
	\item Operations research (1CS-RO);
	\item Machine learning (2CS-ML).
\end{itemize}

\end{frame}


\begin{frame}
\frametitle{Objectives}

\begin{itemize}
	\item Learn some linguistic concepts;
	\item Apply concepts covered in the prerequisite modules to solve NLP tasks;
	\item Discover some NLP tools and resources;
	\item Understand the research directions in this field;
	\item Implement simple solutions to a few problems;
	\item Know how to apply the concepts of language theory to process natural languages;
	\item Get an idea of some of the work carried out within ESI in this area.
\end{itemize}

\end{frame}


\begin{frame}
\frametitle{Course components \& Assessment}

\optword{Course components:}
\begin{itemize}
	\item \textbf{Lectures session:} 1h30/week (can be 2h/week).
	\item \textbf{Practicals session:} 1h30/week.
	\begin{itemize}
		\item \textbf{Labs:} Coding a concept from scratch (graded).
		\item \textbf{Workshops:} Using APIs to solve a problem (guided, not graded).
		\item \textbf{Tutorials:} Exercises (not graded).
	\end{itemize}
	\item \textbf{Quizzes:} Online, after each chapter (graded).
	\item \textbf{Online Demos:} {\scriptsize\url{https://github.com/projeduc/ESI_NLP/tree/main/demos}}.
	\item Some lectures can be presented during practicals' session.
\end{itemize}

\vfill

\optword{Assessment:}
\begin{itemize}
	\item \textbf{Credit:} 3.
	\item \textbf{Continuous assessment:} Labs (20\%) + Quizzes (10\%).
	\item \textbf{Midterm exam:} 30\%.
	\item \textbf{Final exam:} 40\%.
\end{itemize}

\end{frame}

\begin{frame}
\frametitle{Lectures}

\begin{enumerate}
	\item Introduction to NLP.
	\item ML for NLP.
	\item Basic text processing.
	\item Language models.
	\item Part of Speech (PoS) tagging.
	\item Parsing.
	\item Word semantics.
	\item Sentence meaning.
	\item Coreference detection.
	\item Discourse coherence.
	\item Some applications.
\end{enumerate}

\end{frame}

\begin{frame}
	\frametitle{Quizzes}
	
	\begin{itemize}
		\item Answering some questions about the lecture.
		\item \optword{GOAL:} Encouraging course review.
		\item \optword{Method:} Take a Quiz (offline).
		\begin{itemize}
			\item \textbf{When?} Each chapter ending.
			\item \textbf{What?} Google Form.
			\item \textbf{How much?} 10 Quizzes.
			\item \textbf{Duration?} 20 minutes/quiz max.
			\item \textbf{Deadline?} almost a week.
		\end{itemize}
		\item \optword{Assessment:}
		\begin{itemize}
			\item 20 points (the mark is an integer).
			\item In a multiple-choice question, a wrong choice cancels a right one.
			\item For each question, the min mark is 0 (no less).
%			\item la note CC est la moyenne des QCMs.
		\end{itemize}
	\end{itemize}
	
\end{frame}

\begin{frame}
	\frametitle{Labs}
	
	\begin{itemize}
		\item Programming some tasks from scratch.
		\item \optword{GOAL:} Understanding the theoretical notions and the theory/practice transition.
		\item \optword{Languages:} Java, Javascript/nodejs or Python (of your choice).
		\item \optword{Method:} Gaps filling.
		\begin{itemize}
			\item The lab will be implemented in its entirety.
			\item A small part will be removed from the code.
			\item Students must complete this part.
		\end{itemize}
		\item The practical work must be finished face-to-face: 30mn explanation + 1h work.
	\end{itemize}
	
\end{frame}

\begin{frame}
	\frametitle{Tutorials}
	
	\begin{itemize}
		\item Solving some exercises.
		\item \optword{GOAL:} Understanding theoretical notions.
		\item \optword{Method:} Discussion.
		\begin{itemize}
			\item Students try to solve some exercises in groups.
			\item Then, the exercise will be discussed and solved on a whiteboard.
		\end{itemize}
	\end{itemize}
	
\end{frame}

\begin{frame}
	\frametitle{Workshops}
	
	\begin{itemize}
		\item Using APIs to accomplish some tasks.
		\item \optword{GOAL:} Learn how to use existing tools.
		\item \optword{Languages:} Most APIs are in Java or Python.
		\item \optword{Platform:} Jupyter notebook (with Ganymede for Java).
		\item \optword{APIs:} 
		\begin{itemize}
			\item \textbf{Java:} CoreNLP, LangPi, OpenNLP, etc.
			\item \textbf{Python:} NLTK, Spacy, Keras, flair, Gensim, Scikit-learn, tensorflow\_text, tensorflow\_hub, amrlib, etc.
			\item \textbf{Javascript:} jslingua, etc.
		\end{itemize}
		\item \optword{Method:} Discovery
		\begin{itemize}
			\item Guided practical (instruction by instruction); 
			\item Code modifying and observing.
		\end{itemize}
		\item \optword{Online resources:}
		\begin{itemize}
			\item \textbf{URL:} {\scriptsize\url{https://github.com/projeduc/ESI_NLP/tree/main/demos}}
			\item Demos on how to use some APIs.
		\end{itemize}
		\item Not graded.
	\end{itemize}
	
\end{frame}

\begin{frame}
	\frametitle{Color codes}
	
	
	\begin{itemize}
		\item Slides' texts and blocks
		\begin{itemize}
			\item {\Huge\raisebox{-0.125em}{${\color{red}\blacksquare}$}} keywords
			\item {\Huge\raisebox{-0.125em}{${\color{violet}\blacksquare}$}} options (choices)
			\item {\Huge\raisebox{-0.125em}{${\color{olivegreen}\blacksquare}$}} examples
		\end{itemize}
	
		\vfill
		\item Figures (mostly)
		\begin{itemize}
			\item {\Huge\raisebox{-0.125em}{${\color{red}\blacksquare}$}} process
			\item {\Huge\raisebox{-0.125em}{${\color{blue}\blacksquare}$}} inputs/outputs
			\item {\Huge\raisebox{-0.125em}{${\color{green}\blacksquare}$}} intermediate states
		\end{itemize}
	\end{itemize}
	
\end{frame}

\begin{frame}
	\frametitle{Similar courses}
	
	\begin{itemize}
			
			\item Natural Language Processing S21 (updated).
			Carnegie Mellon University. 
			\kurl{http://demo.clab.cs.cmu.edu/NLP/}
			
			\item CS388: Natural Language Processing (2018). 
			University of Texas. 
			\kurl{https://www.cs.utexas.edu/~mooney/cs388/}
			
			\item CSE 517: Natural Language Processing ().
			University of Washington.
			\kurl{https://courses.cs.washington.edu/courses/cse517/}
			
			\item CS 294-5: Statistical Natural Language Processing (2005). 
			University of Berkely. 
			\kurl{https://people.eecs.berkeley.edu/~klein/cs294-5/index.html}
			
			\item CS224n: Natural Language Processing with Deep Learning (updated).
			Stanford university.
			\kurl{https://web.stanford.edu/class/cs224n/}
			
			\item CS 5340/6340: Natural Language Processing (updated). 
			University of Utah.
			\kurl{https://my.eng.utah.edu/~cs5340/schedule.html}
		
	\end{itemize}
	
\end{frame}

\begin{frame}
	\frametitle{Reading resources}
	
	\begin{itemize}
		
		\item \textbf{Traitement automatique du langage naturel - course textbook [FR]}
		
		by Abdelkrime Aries
		
		\kurl{https://github.com/projeduc/ESI_2CS_TALN/releases/download/4.0/taln_v2.pdf}
		
		
		\item \textbf{Speech and Language Processing} 
		
		by Dan Jurafsky and James H. Martin
		
		\kurl{https://web.stanford.edu/~jurafsky/slp3/}
		
		\item \textbf{Natural language processing}
		
		by Jacob Esenstein 
		
		\kurl{https://cseweb.ucsd.edu/~nnakashole/teaching/eisenstein-nov18.pdf}
		
		\item \textbf{Modern Approaches in Natural Language Processing - course textbook (LMU Munich)}
		
		by Carolin Becker, et al.
		
		\kurl{https://slds-lmu.github.io/seminar_nlp_ss20/index.html}
		
		
	\end{itemize}
	
\end{frame}

\begin{frame}
	\frametitle{Office hours}
	
	\begin{itemize}
		
		\item Why?
		\begin{itemize}
			\item Discuss the material being presented in class;
			\item Talk about other related interests (related projects, etc.).
		\end{itemize}
		
		\item Where? 
		\begin{itemize}
			\item New laboratory building (near DG), LCSI lab (first floor), office N 68 (left).
		\end{itemize}
		
		\item When?
		\begin{itemize}
			\item Tuesday, from 13:30 to 16:30.
		\end{itemize}
		
	\end{itemize}
	
\end{frame}

\begin{frame}
	\frametitle{Disclaimer}
	
	\begin{tcolorbox}[colback=red,
		colframe=red,  
		arc=0pt,outer arc=0pt,
		valign=top, 
		halign=center,
		width=\textwidth]
		
		\color{white}
		{\bfseries Disclaimer}
		
	\end{tcolorbox}\vspace{-.5cm}
	\begin{tcolorbox}[colback=my-grey,
		colframe=my-grey,  
		center, arc=0pt,outer arc=0pt,
		valign=top, 
		halign=left,
		width=\textwidth]
		
		\begin{itemize}
			\item An engineer must be a force of good; we must not use our knowledge to harm others.
			\item This course does not encourage malicious use of NLP techniques such as espionage, illegal information gathering, misinformation via chatbots, etc.
			\item Any example that may contain one of these applications is purely educational.	
		\end{itemize}
		
	\end{tcolorbox}
	
\end{frame}


\end{document}

